% Source: https://artofproblemsolving.com/community/c3h2029806_happy_pi_day_maybe_a_competition
% Stable identity: aops-topic-2029806-q09
Let the edge be \(s\). The altitude of a regular tetrahedron is \(s\sqrt{\dfrac{2}{3}}\), so \(s=8\sqrt{\dfrac{3}{2}}=4\sqrt6\). Its volume is
\[
\dfrac{s^3}{6\sqrt2}=\dfrac{(4\sqrt6)^3}{6\sqrt2}=\dfrac{384\sqrt6}{6\sqrt2}=64\sqrt3 .
\]
The inradius is \(\dfrac{1}{4}\) of the altitude: joining the incentre to the four vertices splits the tetrahedron into four congruent pyramids of height \(r\) over the four faces, so \(4\cdot\dfrac{1}{3}Ar=\dfrac{1}{3}A\cdot8\), giving \(r=2\). The figure shows the cross-section through the apex, a base vertex and the midpoint of the opposite base edge.

\begin{asy}
usepackage("amsmath");
size(4.5cm);
defaultpen(fontsize(10pt));
real m=6*sqrt(2); pair V1=(0,0), Md=(m,0), G=(2*m/3,0), Ap=(2*m/3,8);
draw(V1--Md--Ap--cycle,linewidth(0.9)); draw(Ap--G,dashed);
draw(circle((2*m/3,2),2),blue+linewidth(0.9)); dot((2*m/3,2)); draw((2*m/3,2)--(2*m/3+2,2),red); label("$2$",(2*m/3+1,2),N,red);
label("apex",Ap,N); label("$8$",(2*m/3,6),W); label("vertex",V1,S); label("edge midpoint",Md,S);
label("\mbox{}",(min(currentpicture,user=true)+max(currentpicture,user=true))/2,black); // keeps the text after the figure black when LaTeX typesets the labels inline
\end{asy}

The sphere has volume \(\dfrac{4}{3}\pi\cdot2^3=\dfrac{32\pi}{3}\). The space outside the sphere has volume
\[
\boxed{64\sqrt3-\dfrac{32\pi}{3}=\dfrac{192\sqrt3-32\pi}{3}}\approx77.3 .
\]
